\chapter{Concluzii}

\section{Playwright}

Playwright s-a dovedit un instrument robust și flexibil pentru testarea aplicației
NopCommerce pe instanța publică, protejată de Cloudflare. Capacitatea de a intercepta
și valida răspunsuri de rețea prin \texttt{page.expect\_response()} a permis
sincronizarea precisă cu operațiile AJAX ale platformei, eliminând necesitatea
așteptărilor bazate pe timp fix. Mecanismul de auto-wait a redus fragilitatea testelor
în fața variațiilor de latență ale rețelei.

Un avantaj semnificativ a fost integrarea cu ecosistemul Python: fixtures-urile
\texttt{pytest}, tipizarea prin type hints și biblioteca \texttt{playwright-stealth}
au permis construirea unei suite funcționale pe un site cu protecție anti-bot, fără
a recurge la un mediu local. Izolarea contextuală automată între teste a asigurat
că fiecare rulare pornește dintr-o stare curată, fără interferențe de sesiune.

Principala limitare întâmpinată a fost verbozitatea codului în cazul scenariilor
condiționale --- absența soft assertions native a obligat la construcții
\texttt{try/except} repetitive. De asemenea, lipsa unui echivalent vizual al
Cypress App a făcut ca depanarea eșecurilor să fie mai laborioasă, necesitând
activarea Trace Viewer înainte de rulare.

În ansamblu, Playwright reprezintă o alegere solidă pentru suite de testare
complexe, multi-browser, în echipe cu experiență Python sau în scenarii în care
controlul granular al rețelei este esențial.

% -------------------------------------------------------
\section{Cypress}

Cypress a oferit cel mai fluid proces de scriere și depanare a testelor din cadrul
acestui proiect. Funcționalitatea Time Travel Debugging, combinată cu retry-ability-ul
automat, a redus semnificativ timpul necesar identificării și corectării eșecurilor.
API-ul \texttt{cy.intercept()} s-a dovedit intuitiv și expresiv, permițând atât
sincronizarea cu AJAX-ul platformei, cât și inspecția răspunsurilor HTTP cu
sintaxă minimală.

Decizia de a rula testele pe o instanță Docker locală a eliminat problema Cloudflare,
oferind un mediu stabil și reproductibil. TypeScript strict a contribuit la calitatea
codului, detectând erori de selector și de tip la compilare, nu la runtime. Funcțiile
helper tipizate (\texttt{fillAddressForm}, \texttt{selectPaymentMethodAndContinue})
au redus duplicarea codului și au crescut lizibilitatea suitei.

Limitările principale au fost legate de arhitectura framework-ului: suportul exclusiv
pentru Chromium și Firefox, absența unui mecanism de stealth și restricția la un
singur tab de browser au impus compromisuri de infrastructură (Docker) și au limitat
acoperirea cross-browser. Asincronicitatea implicită a \textit{command queue}-ului
Cypress a reprezentat o curbă de învățare pentru scenariile care necesitau accesul
la valori intermediare.

Cypress rămâne alegerea optimă pentru proiecte în care developer experience-ul și
viteza de iterație au prioritate, în special când testele rulează pe o infrastructură
controlată și stack-ul echipei este JavaScript/TypeScript.

% -------------------------------------------------------
\section{De ce se preferă Playwright în industrie}

În ultimii ani, Playwright a câștigat o poziție dominantă în industrie, depășind
Cypress ca framework de referință pentru testarea E2E. Raportul \textit{State of JS
2023} și sondajele comunității indică o adoptare în creștere accelerată, susținută
de mai mulți factori structurali.

\subsection*{Suport multi-browser real}

Playwright testează nativ pe Chromium, Firefox și WebKit cu același cod și aceeași
comandă de rulare. Aceasta este o cerință frecventă în organizațiile care livrează
produse pe multiple platforme și browsere. Cypress acoperă doar Chromium și Firefox,
iar suportul WebKit rămâne experimental --- o limitare semnificativă pentru echipele
care trebuie să garanteze compatibilitatea Safari.

\subsection*{Suport multi-limbaj}

Playwright oferă API-uri oficiale pentru Python, TypeScript, Java și C\#. Această
flexibilitate îl face accesibil echipelor cu stack-uri tehnice diverse și permite
integrarea testelor E2E în proiecte care nu sunt construite pe JavaScript. Cypress
este limitat la JavaScript/TypeScript, ceea ce restrânge adoptarea în organizații
cu echipe mixte.

\subsection*{Arhitectură fără limitări de tab}

Playwright suportă nativ scenarii cu mai multe pagini și tab-uri de browser, ferestre
noi și contexte de browser independente. Cypress, prin arhitectura sa in-browser,
nu poate gestiona mai multe tab-uri simultan --- o limitare care devine vizibilă
în aplicații moderne cu fluxuri OAuth, plăți cu redirect extern sau notificări
în ferestre pop-up.

\subsection*{Viteza de execuție și paralelism}

Playwright rulează testele în paralel pe mai mulți workeri în mod implicit, fără
configurare suplimentară. Izolarea prin \texttt{BrowserContext} este mai ușoară și
mai rapidă decât pornirea unui proces Cypress complet pentru fiecare test. În suite
mari, diferența de timp de execuție devine semnificativă.

\subsection*{Integrare în CI/CD}

Playwright este proiectat explicit pentru medii headless și CI/CD: binarele de browser
sunt descărcate și gestionate automat, nu există dependențe de un display server
(Xvfb), iar rapoartele HTML și fișierele de trace sunt generate nativ. Cypress
necesită configurare suplimentară în CI pentru medii headless și are limitări legate
de licențiere pentru Cypress Cloud (dashboard-ul de raportare).

\subsection*{Tabel comparativ}

\begin{table}[h!]
\centering
\begin{tabular}{lcc}
\toprule
\textbf{Criteriu} & \textbf{Playwright} & \textbf{Cypress} \\
\midrule
Browsere suportate       & Chromium, Firefox, WebKit & Chromium, Firefox \\
Limbaje de scripting     & Python, TS, Java, C\#     & JavaScript / TypeScript \\
Multi-tab / multi-window & Da                        & Nu (limitare arhitecturală) \\
Stealth / bypass bot     & Da (plugin oficial)        & Nu \\
Paralelism nativ         & Da                        & Parțial (Cypress Cloud) \\
Time Travel Debugging    & Trace Viewer (post-mortem) & Da (live, în Cypress App) \\
Developer experience     & Bun                       & Excelent \\
Setup inițial            & Mediu                     & Simplu \\
Integrare CI/CD          & Excelentă                 & Bună \\
\bottomrule
\end{tabular}
\caption{Comparație Playwright vs. Cypress}
\label{tab:comparatie}
\end{table}

\subsection*{Concluzie generală}

Alegerea între Playwright și Cypress depinde de contextul specific al proiectului.
Cypress rămâne superior din punct de vedere al experienței de dezvoltare și al vitezei
de iterație în proiecte mici și medii cu echipe JavaScript. Playwright se impune
în organizații care necesită acoperire cross-browser reală, flexibilitate de limbaj,
scenarii complexe cu mai multe pagini sau integrare nativă în pipeline-uri CI/CD
la scară mare.

Tendința din industrie favorizează Playwright tocmai pentru că aceste cerințe ---
multi-browser, multi-limbaj, CI-first --- sunt din ce în ce mai frecvente în
proiectele moderne, unde calitatea software-ului trebuie garantată pe platforme
și configurații diverse, nu doar pe Chrome.
